\section{Homomorphisms and dimension}
\label{section-homomorphism-dimension}

\noindent
This section contains a collection of easy results relating
dimensions of rings when there are maps between them.

\begin{lemma}
\label{lemma-dimension-going-up}
Suppose $R \to S$ is a ring map satisfying either going up, see
Definition \ref{definition-going-up-down}, or going down
see Definition \ref{definition-going-up-down}.
Assume in addition that $\Spec(S) \to \Spec(R)$
is surjective. Then $\dim(R) \leq \dim(S)$.
\end{lemma}

\begin{proof}
Assume going up.
Take any chain $\mathfrak p_0 \subset \mathfrak p_1 \subset \ldots
\subset \mathfrak p_e$ of prime ideals in $R$.
By surjectivity we may choose a prime $\mathfrak q_0$ mapping
to $\mathfrak p_0$. By going up we may extend this to a chain
of length $e$ of primes $\mathfrak q_i$ lying over
$\mathfrak p_i$. Thus $\dim(S) \geq \dim(R)$.
The case of going down is exactly the same.
See also Topology, Lemma \ref{topology-lemma-dimension-specializations-lift}
for a purely topological version.
\end{proof}

\begin{lemma}
\label{lemma-going-up-maximal-on-top}
Suppose that $R \to S$ is a ring map with the going up property,
see Definition \ref{definition-going-up-down}. If
$\mathfrak q \subset S$ is a maximal ideal.
Then the inverse image of $\mathfrak q$ in $R$
is a maximal ideal too.
\end{lemma}

\begin{proof}
Trivial.
\end{proof}

\begin{lemma}
\label{lemma-integral-dim-up}
Suppose that $R \to S$ is a ring map such that $S$ is integral over $R$.
Then $\dim (R) \geq \dim(S)$, and every closed point of $\Spec(S)$
maps to a closed point of $\Spec(R)$.
\end{lemma}

\begin{proof}
Immediate from Lemmas \ref{lemma-integral-no-inclusion} and
\ref{lemma-going-up-maximal-on-top}
and the definitions.
\end{proof}

\begin{lemma}
\label{lemma-integral-sub-dim-equal}
Suppose $R \subset S$ and $S$ integral over $R$.
Then $\dim(R) = \dim(S)$.
\end{lemma}

\begin{proof}
This is a combination of Lemmas
\ref{lemma-integral-going-up},
\ref{lemma-integral-overring-surjective},
\ref{lemma-dimension-going-up}, and
\ref{lemma-integral-dim-up}.
\end{proof}

\begin{definition}
\label{definition-fibre}
Suppose that $R \to S$ is a ring map.
Let $\mathfrak q \subset S$ be a prime lying
over the prime $\mathfrak p$ of $R$.
The {\it local ring of the fibre at $\mathfrak q$}
is the local ring
$$
S_{\mathfrak q}/\mathfrak pS_{\mathfrak q}
=
(S/\mathfrak pS)_{\mathfrak q}
=
(S \otimes_R \kappa(\mathfrak p))_{\mathfrak q}
$$
\end{definition}

\begin{lemma}
\label{lemma-dimension-base-fibre-total}
Let $R \to S$ be a homomorphism of Noetherian rings.
Let $\mathfrak q \subset S$ be a prime lying
over the prime $\mathfrak p$. Then
$$
\dim(S_{\mathfrak q})
\leq
\dim(R_{\mathfrak p})
+
\dim(S_{\mathfrak q}/\mathfrak pS_{\mathfrak q}).
$$
\end{lemma}

\begin{proof}
We use the characterization of dimension of
Proposition \ref{proposition-dimension}.
Let $x_1, \ldots, x_d$ be elements of $\mathfrak p$
generating an ideal of definition of $R_{\mathfrak p}$ with
$d = \dim(R_{\mathfrak p})$.
Let $y_1, \ldots, y_e$ be elements of $\mathfrak q$
generating an ideal of definition of
$S_{\mathfrak q}/\mathfrak pS_{\mathfrak q}$
with $e = \dim(S_{\mathfrak q}/\mathfrak pS_{\mathfrak q})$.
It is clear that $S_{\mathfrak q}/(x_1, \ldots, x_d, y_1, \ldots, y_e)$
has a nilpotent maximal ideal. Hence
$x_1, \ldots, x_d, y_1, \ldots, y_e$ generate an ideal of definition
of $S_{\mathfrak q}$.
\end{proof}

\begin{lemma}
\label{lemma-dimension-base-fibre-equals-total}
Let $R \to S$ be a homomorphism of Noetherian rings.
Let $\mathfrak q \subset S$ be a prime lying
over the prime $\mathfrak p$. Assume the going down property holds
for $R \to S$ (for example if $R \to S$ is flat, see
Lemma \ref{lemma-flat-going-down}). Then
$$
\dim(S_{\mathfrak q})
=
\dim(R_{\mathfrak p})
+
\dim(S_{\mathfrak q}/\mathfrak pS_{\mathfrak q}).
$$
\end{lemma}

\begin{proof}
By Lemma \ref{lemma-dimension-base-fibre-total}
we have an inequality
$\dim(S_{\mathfrak q}) \leq
\dim(R_{\mathfrak p}) + \dim(S_{\mathfrak q}/\mathfrak pS_{\mathfrak q})$.
To get equality, choose a chain of primes
$\mathfrak pS \subset \mathfrak q_0 \subset \mathfrak q_1 \subset \ldots
\subset \mathfrak q_d = \mathfrak q$ with
$d = \dim(S_{\mathfrak q}/\mathfrak pS_{\mathfrak q})$.
On the other hand, choose a chain of primes
$\mathfrak p_0 \subset \mathfrak p_1 \subset \ldots \subset \mathfrak p_e
= \mathfrak p$ with $e = \dim(R_{\mathfrak p})$.
By the going down theorem we may choose
$\mathfrak q_{-1} \subset \mathfrak q_0$ lying over
$\mathfrak p_{e-1}$. And then we may choose
$\mathfrak q_{-2} \subset \mathfrak q_{-1}$ lying over
$\mathfrak p_{e-2}$. Inductively we keep going until we
get a chain
$\mathfrak q_{-e} \subset \ldots \subset \mathfrak q_d$
of length $e + d$.
\end{proof}

\begin{lemma}
\label{lemma-flat-over-regular-with-regular-fibre}
Let $R \to S$ be a local homomorphism of local Noetherian rings.
Assume
\begin{enumerate}
\item $R$ is regular,
\item $S/\mathfrak m_RS$ is regular, and
\item $R \to S$ is flat.
\end{enumerate}
Then $S$ is regular.
\end{lemma}

\begin{proof}
By Lemma \ref{lemma-dimension-base-fibre-equals-total}
we have
$\dim(S) = \dim(R) + \dim(S/\mathfrak m_RS)$.
Pick generators $x_1, \ldots, x_d \in \mathfrak m_R$ with $d = \dim(R)$,
and pick $y_1, \ldots, y_e \in \mathfrak m_S$
which generate the maximal ideal of $S/\mathfrak m_RS$ with
$e = \dim(S/\mathfrak m_RS)$. Then we see that
$x_1, \ldots, x_d, y_1, \ldots, y_e$ are elements which generate
the maximal ideal of $S$ and $e + d = \dim(S)$.
\end{proof}

\noindent
The lemma below will later be used to show that rings of finite type over
a field are Cohen-Macaulay if and only if they are quasi-finite flat over
a polynomial ring. It is a partial converse to
Lemma \ref{lemma-CM-over-regular-flat}.

\begin{lemma}
\label{lemma-finite-flat-over-regular-CM}
Let $R \to S$ be a local homomorphism of Noetherian local rings.
Assume $R$ Cohen-Macaulay.
If $S$ is finite flat over $R$, or if $S$ is flat over $R$ and
$\dim(S) \leq \dim(R)$, then $S$ is Cohen-Macaulay and $\dim(R) = \dim(S)$.
\end{lemma}

\begin{proof}
Let $x_1, \ldots, x_d \in \mathfrak m_R$ be a regular sequence
of length $d = \dim(R)$. By Lemma \ref{lemma-flat-increases-depth}
this maps to a regular sequence in $S$.
Hence $S$ is Cohen-Macaulay if $\dim(S) \leq d$. This is true
if $S$ is finite flat over $R$ by Lemma \ref{lemma-integral-sub-dim-equal}.
And in the second case we assumed it.
\end{proof}








